\documentclass[a4paper]{article}     
\usepackage{amsfonts}
\usepackage{amsmath}
\usepackage{amssymb}
\usepackage{graphicx}

\newcommand{\Bgcos}{\operatorname{Bgcos}}

\begin{document}

\noindent \large Auteur: Mark Gijsbrechts \\
\noindent \large Studierichting: Informatica\\
\noindent \large Oefeningenreeks 2B nr. 1.8\\

\medskip

\normalsize

Druk $\tan(\Bgcos x)$ uit in $x$:\\

$\tan(\Bgcos x) = \dfrac{\sin(\Bgcos x)}{\cos(\Bgcos x)}$ \\

$\tan(\Bgcos x) = \dfrac{\sin(\Bgcos x)}{x}$ \\ \\

$\alpha = \Bgcos x$ met $\alpha \in \left[ 0, \pi \right]$ \\

$\Leftrightarrow x = \cos \alpha$  \\

$x^2 = \cos^2 \alpha$ \\

$x^2 = 1 - \sin^2 \alpha$ \\

$ \sin^2 \alpha = 1 - x^2$ \\

$ \sin \alpha = \sqrt{1 - x^2}$ \\

$ \sin(\Bgcos x) = \sqrt{1 - x^2}$ \\ \\

$\tan(\Bgcos x) = \dfrac{\sqrt{1 - x^2}}{x}$ \\


\begin{thebibliography}{999}
%\bibitem{a} Referentie
\end{thebibliography}
\end{document} 
